\section{共轭先验让 Bayesian 更新保持同一形式}
\label{sec:13-Machine-Learning/06-Probabilistic-Models/07}

一个落地页改版上线两小时，后台记录到 \(20\) 次访问、\(3\) 次转化。产品经理要一个数字，你手上只有 \(3/20=15\%\)。这个数字的问题不在于算错了，而在于它把二十次访问当成了全部证据，而旧版本三个月来的转化率稳定在 \(20\%\) 附近。更麻烦的是运维要求这块面板每小时刷新一次，你并不想每次都把三个月的原始日志重跑一遍。

Bayesian 更新的公式只有一行，

\[
p(\theta\given x_{1:n})\propto p(\theta)\prod_{i=1}^np(x_i\given\theta),
\]

真正的困难在于右边归一化之后可能掉进一个谁也叫不出名字的分布族，闭式没有，递推也没有。共轭先验挑一个与似然代数结构相配的先验，让后验留在同一族里；上面两个诉求——把历史经验折进估计、只保留一份可增量更新的摘要——就同时被解决了。

\subsection{共轭性来自把同类项收拢}

上一节的形式立刻给出答案。若似然是指数族，

\[
p(x\given\eta)=h(x)\exp\bigl(\eta^{\trans}T(x)-A(\eta)\bigr),
\]

就为自然参数选一个把 \(\eta\) 与 \(A(\eta)\) 都提前写好的先验：

\[
p(\eta\given\chi,\nu)=\exp\bigl(\eta^{\trans}\chi-\nu A(\eta)-B(\chi,\nu)\bigr),
\]

其中 \(B\) 负责对 \(\eta\) 归一化。看到 \(n\) 个独立样本之后，

\[
p(\eta\given x_{1:n},\chi,\nu)
\propto\exp\Bigl[\eta^{\trans}\Bigl(\chi+\sum_{i=1}^nT(x_i)\Bigr)-(\nu+n)A(\eta)\Bigr].
\]

指数里的结构一个字没变，变的只是两个系数：

\[
\chi_n=\chi+\sum_{i=1}^nT(x_i),\qquad \nu_n=\nu+n.
\]

共轭性的来源就这么简单——先验预先摆好了似然会产生的那两种 \(\eta\) 的函数，相乘时只需把系数相加。数据通过充分统计量进入，先验以同样的坐标记录自己的那一份，两者可以直接做加法。

把 \(\nu\) 读成虚拟样本量、\(\chi/\nu\) 读成先验的统计量均值，在很多常见族里说得通，但这是一种解释而非定理：先要保证 \(B(\chi,\nu)\) 有限、先验是正常分布，某些超参数区域里字面解释会失效。

\subsection{Beta--Bernoulli 把成功与失败分开计数}

回到落地页。设 \(X_i\given p\sim\Bernoulli(p)\)，先验取 \(p\sim\operatorname{Beta}(a,b)\)。记 \(s=\sum_iX_i\)，两式相乘并收拢指数，

\[
p\given x_{1:n}\sim\operatorname{Beta}(a+s,\,b+n-s).
\]

成功次数加到第一个参数上，失败次数加到第二个上，别的什么也没发生。用旧版本的经验设 \(a=2,b=8\)，先验均值恰是 \(0.2\)，强度相当于十次虚拟访问。观测到 \(3/20\) 之后，后验是 \(\operatorname{Beta}(5,25)\)，均值 \(5/30\approx0.167\)——不是 \(15\%\)，也不是 \(20\%\)。等样本累积到 \(100\) 次访问、\(12\) 次转化，后验成为 \(\operatorname{Beta}(14,86)\)，均值降到 \(0.14\)，先验的那十次虚拟访问已经被稀释得几乎不起作用。

\begin{center}
\begin{tikzpicture}[>=stealth]

\draw[gray!60,->] (-0.1,0) -- (6.5,0);
\draw[gray!60,->] (0,-0.1) -- (0,3.1);
\foreach \x/\lab in {1.2/{0.1}, 2.4/{0.2}, 3.6/{0.3}, 4.8/{0.4}}{
  \draw[gray!60] (\x,0) -- (\x,-0.1);
  \node[font=\scriptsize,anchor=north] at (\x,-0.12) {\(\lab\)};
}
\node[font=\scriptsize,anchor=north] at (6.3,-0.05) {\(p\)};

\draw[black!45,thick] plot[smooth] coordinates {(0.02,0.03) (0.32,0.37) (0.62,0.59)
 (0.92,0.72) (1.22,0.79) (1.52,0.81) (1.82,0.79) (2.12,0.75) (2.41,0.69) (2.71,0.62)
 (3.01,0.55) (3.31,0.47) (3.61,0.40) (3.91,0.34) (4.21,0.28) (4.51,0.23) (4.80,0.18)
 (5.10,0.14) (5.40,0.11) (5.70,0.09) (6.00,0.06)};

\draw[blue!65,thick] plot[smooth] coordinates {(0.02,0.00) (0.32,0.04) (0.62,0.27)
 (0.92,0.69) (1.22,1.10) (1.52,1.35) (1.82,1.39) (2.12,1.25) (2.41,1.01) (2.71,0.76)
 (3.01,0.52) (3.31,0.34) (3.61,0.21) (3.91,0.12) (4.21,0.06) (4.51,0.03) (4.80,0.02)
 (5.10,0.01) (5.40,0.00) (5.70,0.00) (6.00,0.00)};

\draw[violet,very thick] plot[smooth] coordinates {(0.02,0.00) (0.32,0.00) (0.62,0.03)
 (0.92,0.44) (1.22,1.66) (1.52,2.63) (1.82,2.32) (2.12,1.34) (2.41,0.55) (2.71,0.17)
 (3.01,0.04) (3.31,0.01) (3.61,0.00) (3.91,0.00) (4.21,0.00) (4.51,0.00) (4.80,0.00)
 (5.10,0.00) (5.40,0.00) (5.70,0.00) (6.00,0.00)};

\node[font=\scriptsize,anchor=west,black!55] at (4.3,0.55) {先验 \(\operatorname{Beta}(2,8)\)};
\node[font=\scriptsize,anchor=west,blue!65] at (3.1,1.15) {\(n=20\)：\(\operatorname{Beta}(5,25)\)};
\node[font=\scriptsize,anchor=west,violet] at (2.0,2.75) {\(n=100\)：\(\operatorname{Beta}(14,86)\)};

\end{tikzpicture}
\end{center}

\emph{三条曲线属于同一族，数据只是把两个计数往上加，顺带把宽度收窄。}

图上要看的不是哪条更高，而是三条曲线的形状同属一族：更新没有把分布带到别处去，只是沿着 \((a,b)\) 这两个坐标移动。这正是共轭性买到的东西——后验可以用两个数字完整存下来，下一批数据到达时接着加，历史日志不必保留。后验均值

\[
\E[p\given x]=\frac{a+s}{a+b+n}
\]

同时也是下一次访问转化的预测概率，因为把 \(p\) 积分掉之后 \(P(X_{n+1}=1\given x_{1:n})=\E[p\given x]\)。超参数 \(a,b\) 控制着先验的均值、集中程度与端点行为；\(a<1\) 时得到的是合法的 U 形先验，只是``先验成功次数''这句话在那里失去字面意义。这个族的基本性质见\hyperref[sec:05-Probability/06-Common-Distributions/04]{``Uniform、Beta 与 Dirichlet 分布''}，多类别的推广把 Beta 换成 Dirichlet，第一节 Naive Bayes 的平滑正是它的一次应用。

\subsection{Gamma--Poisson 相加的是暴露量而非样本数}

计数数据是同一套机制的另一个实例。取 \(Y_i\given\lambda\sim\operatorname{Poisson}(\lambda)\)，先验 \(\lambda\sim\operatorname{Gamma}(a,b)\)，密度正比于 \(\lambda^{a-1}e^{-b\lambda}\)。等暴露时长下观测 \(n\) 个样本，

\[
\lambda\given y_{1:n}\sim\operatorname{Gamma}\Bigl(a+\sum_iy_i,\;b+n\Bigr).
\]

若各观测的暴露不同，\(Y_i\sim\operatorname{Poisson}(t_i\lambda)\)，第二个参数更新为 \(b+\sum_it_i\)。与参数相乘的是总暴露量而不是样本条数——一台运行了三百小时的机器和一台运行了三小时的机器，在这个公式里的分量差一百倍。分不清这两者，是计数模型里很常见的一类错误。相关分布的来历见\hyperref[sec:05-Probability/06-Common-Distributions/03]{``Poisson 计数、Exponential 等待与 Gamma 时间''}。

把 \(\lambda\) 积分掉之后的预测分布值得单独看一眼。它不是把后验均值塞回 Poisson 得到的那个 Poisson，而是一个负二项型分布：参数的不确定性额外贡献了一份方差，尾巴比 Poisson 厚。用点估计代入做预测会系统性低估这份波动，样本少时差距很大。

\subsection{Gaussian--Gaussian 相加的是精度}

设观测方差 \(\sigma^2\) 已知，\(X_i\given\mu\sim\mathcal N(\mu,\sigma^2)\)，\(\mu\sim\mathcal N(m_0,\tau_0^2)\)。后验仍是 Gaussian，更新写在精度上：

\[
\tau_n^{-2}=\tau_0^{-2}+n\sigma^{-2},
\qquad
m_n=\tau_n^2\Bigl(\tau_0^{-2}m_0+\sigma^{-2}\sum_ix_i\Bigr).
\]

后验均值是先验均值与样本均值的精度加权平均，权重就是各自贡献的信息量。样本多不自动压倒先验：观测噪声极大时每个样本贡献的精度很小，而一个很窄的先验意味着很强的外部信息或很强的建模承诺。这个更新式与\hyperref[sec:13-Machine-Learning/02-Linear-Models/04]{``Bayesian 线性回归把系数不确定性传到预测''}中的公式是同一个，只是那里的精度换成了矩阵。

方差也未知时，单给 \(\mu\) 一个独立的 Gaussian 先验不再闭合。可用的是 normal--inverse-gamma 一类联合先验，让均值的先验尺度随未知方差一起伸缩。共轭对象常常是整组参数的联合分布，不能逐个参数各挑各的。

\subsection{预测比点估计多走一次积分}

对新观测，

\[
p(x_*\given x_{1:n})=\int p(x_*\given\theta)\,p(\theta\given x_{1:n})\,d\theta,
\]

这一步把后验中仍然说得通的所有参数平均了一遍。共轭模型通常让它闭式可得；非共轭时要靠 Monte Carlo 或变分近似，见\hyperref[sec:13-Machine-Learning/08-Sampling-and-Inference/02]{``MCMC 在无法独立采样时构造平稳链''}与\hyperref[sec:13-Machine-Learning/08-Sampling-and-Inference/03]{``变分推断把后验近似变成优化''}。把 MAP 或后验均值代进似然得到的是 plug-in 分布，它丢掉的正是参数不确定性那一块，样本少时低估得厉害——上面负二项与 Poisson 的差距就是这件事的一个可量化版本，点估计与完整后验的关系另见\hyperref[sec:11-Mathematical-Statistics/02-Point-Estimation/04]{``MAP 把先验与似然合在一个点上''}。

顺序更新是共轭性的一个直接推论：这一批的后验超参数可以原样当作下一批的先验超参数，落地页面板每小时只需存两个数。前提是模型、条件独立结构与参数含义都没有变。数据一旦发生漂移，机械累加旧统计量的后果是模型对一个已经过时的环境越来越自信——曲线越来越窄，指向的却是错的位置。需要遗忘时，可以给旧统计量乘一个小于 \(1\) 的折扣因子，这么做之后它不再对应任何一个干净的 Bayes 模型。

\subsection{计算便利不构成选择先验的理由}

共轭先验对分布的端点、尾部与参数间的相关结构都施加了形状限制，这些限制是被计算便利顺带塞进来的，未必是你想要的。数据很少时超参数会明显影响后验；数据很多但模型错设时，后验只会更集中地支持一个错误的近似——收窄的可信区间在这里是一种误导，讨论见\hyperref[sec:11-Mathematical-Statistics/04-Interval-Estimation/03]{``渐近区间与 Bayesian 可信区间''}。检查先验是否荒谬有一个便宜办法：从先验里采参数，再从似然生成数据，看生成的数据在量级和结构上像不像真实业务能产生的东西。落地页的例子里，若先验采样生成出一批转化率高于 \(80\%\) 的场景，那个先验就该换掉。

超参数本身还有三种处理方式，它们的语义并不相同。固定超参数表达的是分析之前就持有的外部假设，最诚实但要求你说得出理由。Empirical Bayes 用同一批数据把超参数估出来，方便且常常表现不错，但这一步的不确定性没有被传播下去，后验因此偏窄。层次 Bayes 再给超参数一层先验，语义完整，牺牲掉的是闭式解：它一般就此消失，得回到采样或变分。

共轭性最值得留下的东西，不是那几张速查表，而是它把``数据提供了哪些充分统计量、先验提供了多少同量纲的信息''摆在同一组坐标里，两边加起来就是后验。下次遇到非共轭模型时，这套结构仍能帮你定位困难：是似然的形状不允许收拢同类项，是参数受了约束，还是归一化常数根本算不出来。
